\begin{frame}[fragile]{손으로 한 번: kNN 예측 직접 계산}
  학습 데이터 6개, 두 반(불합격 0 / 합격 1):
  \[(1,1)\to0,\ (2,1)\to0,\ (1,2)\to0,\ (8,8)\to1,\ (9,8)\to1,\ (8,9)\to1\]
  새 점 $(7,7)$을 $k=3$으로 분류 --- 유클리드 거리 계산:
  \small
  \begin{tabular}{@{}lll@{}}
    \toprule
    학습 데이터 & 라벨 & $(7,7)$까지 거리 \\
    \midrule
    (8,8) & 1 & $\sqrt{2}\approx1.414$ \\
    (9,8) & 1 & $\sqrt{5}\approx2.236$ \\
    (8,9) & 1 & $\sqrt{5}\approx2.236$ \\
    (2,1) & 0 & $\sqrt{61}\approx7.810$ \\
    (1,2) & 0 & $\sqrt{61}\approx7.810$ \\
    (1,1) & 0 & $\sqrt{72}\approx8.485$ \\
    \bottomrule
  \end{tabular}
  \vskip0.8em
  가장 가까운 3개 --- (8,8), (9,8), (8,9) --- 가 \textbf{모두
  라벨 1} $\Rightarrow$ 다수결로 $(7,7)$은 \kb{1(합격)}.
  \vskip0.4em
  \texttt{knn\_predict}가 정확히 이 정렬·다수결 과정을 자동화한 것.
  \textbf{회귀}로 보면: 같은 라벨을 $1,1,1,4,4,4$로 두면 가장 가까운
  3개 라벨의 평균 $(4+4+4)/3 = 4.0$ --- \textbf{같은 $k$개의
  이웃}을 쓰고, 집계 방식(다수결 vs 평균)만 다르다.
\end{frame}
